% Part: first-order-logic
% Chapter: axiomatic-deduction
% Section: deduction-theorem-quantifiers

\documentclass[../../../include/open-logic-section]{subfiles}

\begin{document}

\olfileid{fol}{axd}{ddq}

\olsection{مبرهنة الاستنباط مع المكمّمات}

\begin{thm}[مبرهنة الاستنباط]
\ollabel{thm:deduction-thm-q} تكون~${\OLId{greek-variable}{Gamma}} \cup \{!A\} \Proves !B$
إذا وفقط إذا كان~${\OLId{greek-variable}{Gamma}} \Proves !A \lif !B$.
\end{thm}

\begin{proof}
% استكمال برهاني (0120-I1): أُثبت الاتجاه العكسي، وأُتمّت حالة QR الوجودية التي تركها الأصل الإنجليزي تمرينًا.
أما اتجاه ``إذا'' فهو الاتجاه اليسير المثبت في
\olref[ded]{thm:deduction-thm}، ولا يتوقف على صورة آخر قاعدة. وأما اتجاه
``فقط إذا'' للاشتقاقات الكمية فنتمّ الاستقراء السابق بإضافة صورتي~\QR{}.

نعيد الاستقراء على طول !!{derivation} الصيغة~$!B$
من~${\OLId{greek-variable}{Gamma}} \cup \{!A\}$.

أما الأساس فبرهانه بعينه برهان الأساس في
\olref[ded]{thm:deduction-thm}.

وأما الخطوة، فنفرض أن !!{derivation} الصيغة~$!B$ من
${\OLId{greek-variable}{Gamma}} \cup \{!A\}$ ينتهي إلى~$!B$ بقاعدة استدلال.
إن كانت إثبات المقدَّم، جرى برهان
\olref[ded]{thm:deduction-thm}. وإن كانت~\QR، فنأخذ أولًا الحالة
$!B \ident !C \lif \lforall[{\OLId{latin-ordinary}{x}}][!D({\OLId{latin-ordinary}{x}})]$.
وقد سبقت !!a{formula} صورتها~$!C \lif !D({\OLId{latin-ordinary}{a}})$
في ال!!{derivation}، ولم يرد~${\OLId{latin-ordinary}{a}}$ في~$!C$ ولا~$!A$ ولا~${\OLId{greek-variable}{Gamma}}$.
فيكون لدينا
\begin{align*}
  {\OLId{greek-variable}{Gamma}} \cup \{!A\} & \Proves !C \lif !D({\OLId{latin-ordinary}{a}}),\\
  \intertext{فيشمله فرض الاستقراء، ونحصل على}
    {\OLId{greek-variable}{Gamma}} & \Proves !A \lif (!C \lif !D({\OLId{latin-ordinary}{a}})).\\
  \intertext{ثم نأخذ}
  & \Proves (!A \lif (!C \lif !D({\OLId{latin-ordinary}{a}}))) \lif ((!A \land !C) \lif !D({\OLId{latin-ordinary}{a}}))\\
  \intertext{ونجري إثبات المقدَّم فينتج}
  {\OLId{greek-variable}{Gamma}} & \Proves (!A \land !C) \lif !D({\OLId{latin-ordinary}{a}}).\\
  \intertext{ولم يختل شرط المتغير المميَّز، فيصح إلحاق خطوة بهذا ال!!{derivation} بقاعدة \QR، فيحصل}
    {\OLId{greek-variable}{Gamma}} & \Proves (!A \land !C) \lif \lforall[{\OLId{latin-ordinary}{x}}][!D({\OLId{latin-ordinary}{x}})].\\
    \intertext{ولدينا كذلك}
    & \Proves ((!A \land !C) \lif \lforall[{\OLId{latin-ordinary}{x}}][!D({\OLId{latin-ordinary}{x}})]) \lif (!A \lif (!C \lif \lforall[{\OLId{latin-ordinary}{x}}][!D({\OLId{latin-ordinary}{x}})])),\\
    \intertext{فينتهي إثبات المقدَّم إلى}
    {\OLId{greek-variable}{Gamma}} & \Proves !A \lif (!C \lif \lforall[{\OLId{latin-ordinary}{x}}][!D({\OLId{latin-ordinary}{x}})]),
\end{align*}
أي إن ${\OLId{greek-variable}{Gamma}} \Proves !A \lif !B$.

وأما صورة~\QR{} الأخرى، فإذا كان
$!B \ident \lexists[{\OLId{latin-ordinary}{x}}][!D({\OLId{latin-ordinary}{x}})] \lif !C$، فقد سبقت في ال!!{derivation}
الصيغة~$!D({\OLId{latin-ordinary}{a}}) \lif !C$، ولم يرد~${\OLId{latin-ordinary}{a}}$ في~$!C$ ولا~$!A$ ولا~${\OLId{greek-variable}{Gamma}}$.
فيعطينا فرض الاستقراء
\[
  {\OLId{greek-variable}{Gamma}} \Proves !A \lif (!D({\OLId{latin-ordinary}{a}}) \lif !C).
\]
ونبدّل ترتيب المقدِّمين بقضية صادقة منطقيًا، فينتج
\[
  {\OLId{greek-variable}{Gamma}} \Proves !D({\OLId{latin-ordinary}{a}}) \lif (!A \lif !C).
\]
وشرط الثابت المميَّز باقٍ، فتعطينا~\QR{}
\[
  {\OLId{greek-variable}{Gamma}} \Proves \lexists[{\OLId{latin-ordinary}{x}}][!D({\OLId{latin-ordinary}{x}})] \lif (!A \lif !C).
\]
ثم نعيد ترتيب المقدِّمين بقضية صادقة منطقيًا أخرى، ونجري إثبات المقدَّم،
فنخلص إلى
\[
  {\OLId{greek-variable}{Gamma}} \Proves !A \lif (\lexists[{\OLId{latin-ordinary}{x}}][!D({\OLId{latin-ordinary}{x}})] \lif !C),
\]
أي~${\OLId{greek-variable}{Gamma}} \Proves !A \lif !B$. فبهاتين الصورتين لـ~\QR{}، وبحالات
الفرض والبديهية وإثبات المقدَّم المتقدمة، يتم الاستقراء.
\end{proof}

\begin{prob}
  تحقّق تفصيلًا من بقاء شروط الثابت المميَّز في صورتي~\QR{} في برهان
  \olref[fol][axd][ddq]{thm:deduction-thm-q}.
\end{prob}

\end{document}
